\section{Алгоритмы информационно-поисковой системы}
\label{sec:algorithms}

\subsection{Предобработка текста}

Исходный текст проходит следующую последовательность преобразований.
Параметр \texttt{lang} определяет язык предобработки; он устанавливается
автоматически при загрузке документа (классификатор записывает результат
в колонку \texttt{lang} таблицы \texttt{documents}) и при поиске определяется
для запроса тем же классификатором.

\begin{enumerate}
	\item приведение к нижнему регистру;
	\item удаление цифр, пунктуации и спецсимволов; для французского режима
	сохраняются диакритические символы (\texttt{àâäéèêëîïôöùûüç\"{o}e});
	\item токенизация по пробелам;
	\item фильтрация стоп-слов (соответствующий набор для каждого языка
	из NLTK) и односимвольных токенов;
	\item стемминг алгоритмом Портера (только для английского; французский
	текст не стемминуется, чтобы сохранить окончания).
\end{enumerate}

Модуль \texttt{lang\_text.py} реализует отдельный тракт предобработки
с сохранением диакритики для классификации языка; поисковый тракт
(\texttt{document\_processor.py}) работает в зависимости от параметра
\texttt{lang}.

Результат предобработки "--- список основ терминов документа.

\subsection{Построение словаря и вычисление IDF}

Словарь строится по всему корпусу документов и представляет собой упорядоченное
множество уникальных терминов. Индекс термина в словаре определяет позицию
соответствующей компоненты в векторе документа.

Инверсная документная частота каждого термина вычисляется по формуле:

\begin{equation}
	\operatorname{IDF}_i = \log\left(\frac{N}{n_i + 1}\right) + 1,
	\label{eq:idf_full}
\end{equation}

где $N$ "--- общее число документов в коллекции; $n_i$ "--- число документов,
содержащих термин $i$. Слагаемое $+1$ в знаменателе и в результате гарантирует
положительные значения IDF для терминов, встречающихся во всех документах.
Добавление $1$ к результату гарантирует, что термины, присутствующие в каждом
документе, получат ненулевой вес, а не нулевой.

\subsection{TF-IDF-векторизация}

Вес термина $i$ в документе $j$ вычисляется как произведение:

\begin{equation}
	w^{j}_{i} = \operatorname{TF}^{j}_{i} \times \operatorname{IDF}_i,
	\label{eq:tfidf_full}
\end{equation}

где $\operatorname{TF}^{j}_{i} = 1 + \log(\text{tf}^{j}_{i})$ при
$\text{tf}^{j}_{i} > 0$ "--- логарифмическая нормализация частоты термина,
снижающая влияние часто повторяющихся слов. При нулевой частоте TF принимает
значение $0$.

Полученный вектор нормируется евклидовой нормой:

\begin{equation}
	\vec{v} = \frac{\vec{w}}{\lVert \vec{w} \rVert},
	\qquad
	\lVert \vec{w} \rVert = \sqrt{\sum_{k} w_k^2}.
	\label{eq:vector_norm}
\end{equation}

Нормировка обеспечивает независимость оценки близости от длины документа:
косинусная близость между двумя нормированными векторами равна их скалярному
произведению.

Векторизация выполняется функцией \texttt{vectorize\_text\_with\_dim(text, dim)}:
размерность \texttt{dim} фиксирована и равна 5000. Если размер словаря превышает
размерность, компоненты с индексами $\ge dim$ отбрасываются.

\subsection{Векторный поиск}

Обработка поискового запроса:

\begin{enumerate}
	\item запрос предобрабатывается аналогично документам (токенизация, стемминг,
	      удаление стоп-слов);
	\item запрос векторизуется той же функцией \texttt{vectorize\_text\_with\_dim};
	\item в PostgreSQL выполняется поиск оператором \texttt{<=>} (косинусное расстояние):
\end{enumerate}

\begin{lstlisting}[caption={Запрос векторного поиска в БД.},label={lst:vector_search}]
SELECT id, title, content,
       1 - (embedding <=> %s::vector) AS similarity
FROM documents
ORDER BY embedding <=> %s::vector
LIMIT %s;
\end{lstlisting}

Релевантность документа запросу вычисляется как косинусная близость:

\begin{equation}
	\operatorname{sim}(d, q) = \frac{(\vec{D}, \vec{Q})}{\lVert \vec{D} \rVert \cdot \lVert \vec{Q} \rVert}
	= 1 - d_{\cos}(\vec{D}, \vec{Q}),
	\label{eq:cosine_full}
\end{equation}

где $d_{\cos} = \vec{D} \times \vec{Q}^{\top}$ "--- косинусное расстояние
(возвращается оператором \texttt{<=>}). Результаты с нулевой близостью
отбрасываются.

\subsection{Формирование сниппета с подсветкой}

Для каждого найденного документа строится сниппет "--- фрагмент текста длиной
до 500 символов с подсветкой релевантных терминов:

\begin{enumerate}
	\item в тексте документа находится первое вхождение любого термина запроса
	      (с учётом стемминга);
	\item выделяется окно вокруг найденной позиции (по 160 символов в каждую
	      сторону, ограниченное границами документа);
	\item все вхождения терминов запроса в выделенном фрагменте оборачиваются
	      в HTML-тег \texttt{<mark>};
	\item если ни один термин запроса не встретился в документе, сниппет строится
	      по ключевым терминам самого документа (топ-5 по TF-IDF).
\end{enumerate}

Подсветка применяется как в сниппетах поисковой выдачи, так и при просмотре
полного текста документа.

\subsection{Индексирование нового документа}

При загрузке документа (пользовательский ввод или файл) выполняется:

\begin{enumerate}
	\item извлечение текста (прямое для \texttt{.txt}, через библиотеки
	      \texttt{pypdf} / \texttt{python-docx} для соответствующих форматов,
	      очистка HTML/RTF);
	\item предобработка текста (токенизация, стемминг, фильтрация);
	\item пересчёт словаря и IDF по всему текущему корпусу + новый документ;
	\item сохранение обновлённых словаря и IDF в БД;
	\item вычисление TF-IDF-вектора нового документа;
	\item вставка документа с вектором в таблицу \texttt{documents}.
\end{enumerate}

Словарь и IDF хранятся в базе данных (таблицы \texttt{vocabulary} и
\texttt{idf\_values}) для персистентности, а также дублируются в памяти
(глобальные структуры \texttt{VOCAB} и \texttt{IDF} в модуле
\texttt{document\_processor.py}) для исключения обращений к БД при каждом
поисковом запросе.

\subsection{Вычисление метрик качества поиска}

Для оценки качества используются стандартные информационно-поисковые метрики.
Пусть $Rel$ "--- множество релевантных документов по запросу, $R$ "--- множество
выданных системой документов.

\begin{itemize}
	\item \textbf{Точность (Precision):}
	\begin{equation}
		P = \frac{|R \cap Rel|}{|R|}.
		\label{eq:precision}
	\end{equation}
	\item \textbf{Полнота (Recall):}
	\begin{equation}
		R_{\text{rec}} = \frac{|R \cap Rel|}{|Rel|}.
		\label{eq:recall}
	\end{equation}
	\item \textbf{F-мера (F-score):}
	\begin{equation}
		F_{\beta} = \frac{(1 + \beta^{2}) \cdot P \cdot R_{\text{rec}}}
		                {\beta^{2} \cdot P + R_{\text{rec}}}.
		\label{eq:fscore}
	\end{equation}
\end{itemize}

Параметр $\beta$ регулирует баланс: при $\beta = 1$ точность и полнота равнозначны,
при $\beta < 1$ приоритет у точности, при $\beta > 1$ "--- у полноты.
В системе вычисляются $F_1$, $F_{0.5}$ и $F_2$.

\subsection{Миграции базы данных}

Схема базы данных управляется собственным механизмом миграций (модуль
\texttt{migrate.py}). Миграции представляют собой пронумерованные SQL-файлы
из каталога \texttt{backend/migrations/}. Порядок выполнения отслеживается в
таблице \texttt{schema\_migrations}: перед выполнением проверяется,
не была ли миграция уже применена.

\begin{itemize}
	\item \texttt{001\_init.sql} "--- создаёт расширение \texttt{pgvector},
	      таблицы \texttt{documents} (с полем \texttt{embedding vector(5000)}),
	      \texttt{search\_logs} и \texttt{schema\_migrations};
	\item \texttt{002\_vocab\_idf.sql} "--- создаёт таблицы \texttt{vocabulary}
	      и \texttt{idf\_values} для персистентного хранения словаря и IDF;
	\item \texttt{003\_file\_watchers.sql} "--- добавляет столбцы \texttt{summary}
	      и \texttt{file\_path} в таблицу \texttt{documents}, создаёт таблицы
	      \texttt{watch\_clients} и \texttt{watch\_events} для поддержки
	      файлового наблюдателя.
\end{itemize}
